\section{Conclusions}\label{sec:conclusions}

\subsection{Conclusions}\label{subsec:conclusion-summary}

This thesis set out to design, implement, and evaluate a controlled proof-of-concept for workflow-based IoT-style action orchestration. The implemented system combined PC-hosted n8n, Python middleware, shared JSON contracts, a documented validation design, CSV-based evaluation, and Raspberry Pi middleware validation. The result is a focused implementation that demonstrates a complete temperature-to-fan event-action path while keeping the scope measurable and transparent.

\textbf{RQ1: How accurately and consistently does the deterministic n8n workflow produce the expected fan action across the defined temperature test cases?}
The deterministic workflow produced the expected fan actions for the defined high-temperature and low-temperature cases. The high-temperature case led to \texttt{fan\_on}, and the low-temperature case led to \texttt{fan\_off}. This behavior was predictable because the workflow used an explicit threshold rule. The conclusion is limited to the defined test cases and the small run count reported in Chapter~\ref{sec:results}; it should not be interpreted as evidence of broad performance across many sensor conditions.

\textbf{RQ2: How accurately and consistently does the minimal agent-enhanced workflow produce contract-valid action output compared with the deterministic baseline?}
The minimal agent-enhanced workflow produced structured, contract-shaped action output for the defined cases and could be routed through the same action pipeline as the deterministic baseline. This shows that an LLM-based decision step can be placed behind the same contract-oriented boundary used by the deterministic workflow in this proof-of-concept. The result does not establish broad agent autonomy, model benchmarking, or a multi-agent architecture. It shows contract-shaped output behavior for the implemented cases.

\textbf{RQ3: What latency, resource, and deployment behavior is observed during local execution and Raspberry Pi middleware validation?}
The CSV-based evaluation captured action outcomes and latency evidence for the local workflows. The Raspberry Pi middleware validation showed that PC-hosted n8n could call the Pi-hosted middleware/action endpoint over HTTP and receive the expected simulated fan action responses. Resource and thermal observations were also collected as supporting evidence. These observations are useful for validating the defined deployment pattern, but they are not a large-scale benchmark and do not represent a full n8n deployment on Raspberry Pi.

Overall, the thesis demonstrates a small but complete workflow-to-action architecture. The implementation did not include full n8n-on-Raspberry-Pi deployment, real GPIO sensor or fan hardware, physical fan control, production-grade safety enforcement, a broad multi-agent architecture, or large-scale benchmarking. Within the defined scope, however, it shows how workflow automation, middleware, structured contracts, validation design, and measurement can be combined into a coherent action pipeline.

\subsection{Main Contributions}\label{subsec:main-contributions}

The main contribution of the thesis is a complete temperature-to-fan event-action proof-of-concept that connects workflow orchestration to action execution through an explicit middleware boundary. The system includes a deterministic baseline workflow, a minimal agent-enhanced workflow, shared JSON contracts, a documented validation design, and a lightweight CSV-based evaluation harness.

A second contribution is the Raspberry Pi middleware validation setup. By keeping n8n on the PC and running the Python middleware/action endpoint on separate Raspberry Pi hardware, the thesis validates the networked edge-side part of the architecture without turning the work into a full edge deployment project. This supports the broader argument that even a small IoT-style action pipeline benefits from clear boundaries, structured action messages, and measurable evidence.

\subsection{Future Work}\label{subsec:future-work}

Future work could extend the proof-of-concept by deploying the full n8n stack on Raspberry Pi or another edge device and comparing that deployment with the PC-hosted workflow used in this thesis. Another natural extension would be to replace the simulated \texttt{fan\_on} and \texttt{fan\_off} endpoints with GPIO-connected sensor and actuator hardware, while preserving the middleware and contract boundary.

The evaluation could also be expanded with more temperature cases, repeated runs, and clearer visual summaries of latency and resource behavior. The current validation design could be developed into stronger runtime safety enforcement, including stricter schema validation, clearer approval handling, and more robust rejection of unknown or risky actions. The agent-enhanced workflow could be evaluated with additional prompts, models, and memory strategies, provided that such work remains grounded in measurable action outputs rather than broad claims about autonomy.

Finally, the architecture could be adapted to more realistic IoT contexts, such as smart building control, industrial monitoring, environmental sensing, or health and environment monitoring. These extensions would need their own validation and safety analysis, but they follow naturally from the same pattern used here: receive an event, convert it into a structured action decision, validate the action boundary, execute through middleware, and record evidence.

\subsection{Final Remarks}\label{subsec:final-remarks}

The value of this thesis is not the fan action itself. The fan scenario is a controlled proxy for a larger engineering problem: how workflow automation and minimal agent-enhanced decision making can be connected to action execution in a way that remains structured, observable, and bounded. The thesis shows that a focused workflow-to-action architecture can be implemented and evaluated with modest components, while still making the limitations of the system visible.
